\section{Marco Teórico}
\label{sec:marco-teorico}

\subsection{Repositorios Institucionales}

Un repositorio institucional (RI) es un archivo en línea donde se depositan, en formato digital, materiales derivados de la producción científica o académica de una institución. Los repositorios institucionales se han convertido en la principal forma de publicar, preservar y difundir la información digital de las organizaciones \textcite{wikipedia2024repositorio}. En el ámbito universitario, los RI permiten centralizar, visibilizar y dar acceso a la producción intelectual generada por estudiantes, docentes e investigadores, constituyéndose en un pilar fundamental para la memoria institucional.

\subsection{Scrum}

Scrum es un marco de trabajo ágil para la gestión y desarrollo de proyectos, caracterizado por su enfoque iterativo e incremental. \textcite{pressman2014} describe Scrum como un proceso que enfatiza la colaboración del equipo, la entrega continua de valor y la adaptación a los cambios mediante ciclos cortos denominados sprints. Scrum define roles específicos (Product Owner, Scrum Master, equipo de desarrollo), artefactos (product backlog, sprint backlog, incremento) y ceremonias (Sprint Planning, Daily Standup, Sprint Review, Sprint Retrospective) que facilitan la organización y el control del proyecto.

\subsection{UWE (UML-based Web Engineering)}

UWE es una metodología de desarrollo de aplicaciones web basada en UML (Unified Modeling Language), propuesta por \textcite{koch2001uwe}. UWE extiende UML con estereotipos específicos para el modelado de sistemas web, permitiendo representar de manera sistemática los modelos de contenido, navegación, presentación y procesos. Esta metodología proporciona un enfoque estructurado para el análisis y diseño de aplicaciones web, facilitando la transición desde los requisitos hasta la implementación.

\subsection{Trazabilidad Documental}

La trazabilidad documental es la capacidad de rastrear el ciclo de vida de un documento a través de las distintas etapas por las que atraviesa, desde su creación hasta su versión final. En el contexto académico, la trazabilidad permite conocer en todo momento el estado de un proyecto de tesis o trabajo de grado, quién lo ha revisado, qué observaciones se han realizado y qué decisiones se han tomado, garantizando la transparencia y el control del proceso.

\subsection{Tecnologías del Sistema}

\subsubsection{Laravel y PHP}

Laravel es un framework de desarrollo web basado en PHP que sigue el patrón de arquitectura MVC (Modelo-Vista-Controlador). Proporciona un conjunto de herramientas y componentes que facilitan tareas comunes como el enrutamiento, la autenticación, el acceso a bases de datos y la generación de interfaces \textcite{laravel2024docs}. PHP es un lenguaje de programación del lado del servidor ampliamente utilizado en el desarrollo de aplicaciones web dinámicas.

\subsubsection{PostgreSQL}

PostgreSQL es un sistema de gestión de bases de datos relacional y orientado a objetos, de código abierto. Destaca por su robustez, cumplimiento de estándares SQL, soporte para transacciones ACID y extensibilidad \textcite{postgresql2024docs}. Es adecuado para aplicaciones que requieren integridad de datos y consultas complejas, como el sistema de gestión académica propuesto.

\subsubsection{React e Inertia}

React es una biblioteca de JavaScript para la construcción de interfaces de usuario basadas en componentes reutilizables \textcite{react2024docs}. Inertia.js actúa como un adaptador que permite construir aplicaciones de una sola página (SPA) utilizando el enrutamiento y los controladores del lado del servidor, sin necesidad de construir una API REST independiente \textcite{inertia2024docs}. La combinación de Laravel e Inertia con React permite desarrollar aplicaciones modernas con una experiencia de usuario fluida.

\subsection{Investigaciones Relacionadas}

\subsubsection{Repositorios y Acceso Abierto en Bolivia}

\textcite{perales2022repositorios} analizaron el desarrollo de repositorios y revistas científicas de acceso abierto en Bolivia, evidenciando los avances y rezagos en la comunicación científica universitaria. El estudio compiló información de más de medio centenar de portales de internet, concluyendo que, si bien existen iniciativas de acceso abierto, aún persisten limitaciones en la implementación y sostenibilidad de los repositorios institucionales.

\subsubsection{Repositorio Institucional en la Universidad de Sevilla}

\textcite{leon2012implantacion} presentaron la implantación de un repositorio de contenidos institucional en la Universidad de Sevilla, basado en tres ejes fundamentales: recolección, preservación y divulgación de la información. El artículo detalla los aspectos de definición, planificación y desarrollo del proyecto, así como los requisitos técnicos y de gestión considerados.

\subsubsection{Acceso Abierto y el Repositorio COLIBRI}

\textcite{seroubian2022acceso} aborda el movimiento de ciencia abierta y su relación con el acceso abierto, presentando la experiencia del repositorio institucional COLIBRI de la Universidad de la República (Uruguay). El trabajo señala los avances en torno a la política de acceso abierto en el país y la repercusión en el repositorio, destacando la importancia de los RI como pilares de la comunicación científica.

\subsubsection{PROYECTOSAPP: Sistema Web para la Gestión de Repositorio Digital}

\textcite{mesa2019proyectosapp} desarrollaron un sistema web para la gestión de repositorio digital de proyectos de investigación en la Universidad de San Buenaventura (Colombia). El proyecto, desarrollado bajo la metodología XP, abordó problemas similares a los identificados en la FICH: ausencia de un sistema que muestre, organice y preserve los trabajos y proyectos presentados por estudiantes y docentes. Los resultados incluyeron diagramas de procesos BPMN, requerimientos funcionales y no funcionales, y diagramas de casos de uso UML.

\endinput
